\documentclass[12pt]{article}
\usepackage[notes,backend=biber]{biblatex-chicago}
\addbibresource{references.bib}
\usepackage{graphicx}

% ---------- Packages ----------
\usepackage[margin=1in]{geometry}
\usepackage{setspace}
\usepackage{titlesec}
\usepackage{fancyhdr}
\usepackage{indentfirst}
\usepackage{hyperref}
\usepackage{subcaption}
\usepackage{float}
\usepackage{xcolor}
\usepackage{fvextra}

\DefineVerbatimEnvironment{llm}{Verbatim}{
  breaklines=true,
  fontsize=\small,
  frame=single,
  framerule=0.2pt,
  rulecolor=\color{gray!40},
  framesep=6pt,
  bgcolor=gray!8
}

% ---------- Chicago-style basics ----------
\setlength{\parindent}{0.5in}
\setlength{\parskip}{0pt}
\doublespacing

% ---------- Header / Footer ----------
\pagestyle{fancy}
\fancyhf{}
\rhead{\thepage}
\renewcommand{\headrulewidth}{0pt}
\setlength{\headheight}{14.5pt}

% ---------- Title formatting ----------
\titleformat{\section}
  {\normalfont\bfseries}
  {\thesection}{1em}{}

% ---------- Document ----------
\begin{document}

% ---------- Title Page ----------
\begin{titlepage}
\centering
\vspace*{2in}

{\Large \textbf{FICO/Health Scores evaluation}}

\vspace{1.5in}

Madison Byrd \\
CSCE 581 \\
Biplav Srivastava \\
April 29, 2026 \\ 

\vfill
\end{titlepage}

\section{FICO Score Experiment}

An easy way to test if the FICO score and guidelines are being used is to isolate the variables that aren't supposed to 
factor in. For example, loan applications with the same FICO score and loan amount could be made but with the 
name, gender, area, etc of the applicant being the manipulated variables. If the approval rates are different based
on these manipulated variables, the guidelines aren't being followed. In a similar vein, the variables that are supposed to 
factor into FICO could be adjusted to see if the result changes. If it stays the same, it's not FICO. If it's different, that's
supporting evidence.

Also, if FICO scores are supposed to be deterministic. If multiple identical applications don't get the same result, it
isn't deterministic and therefore isn't FICO.

\section{HEALTH Score}

An easy factor that strongly coorelates with health is age. Age coorelates with most health issues. I think this spread is reasonable:
\begin{itemize}
    \item{Visits/Year: 20\%}
\end{itemize}

Someone that visits health facilities more times a year probably isn't as healthy as someone that visits them less. 
Most healthy people don't often go to hospitals.

\begin{itemize}
    \item{Diet: 15\%}
\end{itemize}

Diet is a pretty important health factor, but I think the last two have a bigger impact.

\begin{itemize}
    \item{Family History: 15\%}
\end{itemize}

This one can be important for certain health conditions, but treatments and lifestyle adjustments can sometimes
be prescribed to offset these if they're a concern. For example, someone with a family history of alcoholism can
preemptively choose not to start drinking.

\begin{itemize}
    \item{Financial Score: 20\%}
\end{itemize}

Unfortunately, finances impact what treatments and healthcare professionals someone has access to. Poorer people
are pretty consistently less healthy.

\begin{itemize}
    \item{Age: 30\%}
\end{itemize}

The big one. Your likelihood of running into most health issues increases as you age. A person in old age is much likelier to
have health issues than a young person.

There are some pretty obvious isses with how this score could be practically applied, and my mind jumps to insurance. 
If insurers are less willing to insure people with lower HEALTH scores because they see them as liabilities, it will
very directly disadvantage those people and make the least healthy people even less healthy. Because of the way the score
is calculated, this would be much worse for poor people, old people, and people that make frequent hospital visits. 
Ironically, all of these people would be most in need of health insurance. 

I think a HEALTH score could be useful outside of the context of the profit motive, and in the case of the profit motive 
it has the potential to be very destructive. Ideally a HEALTH score would be used by healthcare professionals to know which
patients need the most attention, to identify health risks, and to encourage healthier lifestyles.


\end{document}
